\section{相关技术基础}

实现一个用户态文件系统至少要回答四个问题：内核与用户空间如何交互？
元数据如何持久化？文件操作应满足哪些语义约束？并发修改怎样保证一致性？
RucksFS 的选择分别是 FUSE 框架、RocksDB 存储引擎、POSIX 语义规范和悲观并发控制。

\subsection{FUSE 用户态文件系统框架}

传统文件系统（如 ext4、XFS、Btrfs）以内核模块的形式运行，
开发和调试需要修改内核源码、编译模块、重启系统，迭代周期长且内核 bug 可能导致整机崩溃。
FUSE（Filesystem in Userspace）\cite{szeredi2005fuse} 提供了一种将文件系统逻辑搬到用户空间的替代方案。
自 Linux 2.6.14 起，FUSE 被纳入内核主线，成为用户态文件系统开发的事实标准。
JuiceFS\cite{juicefs2021}、CubeFS\cite{liu2019cubefs}、GlusterFS、SSHFS 等主流分布式或网络文件系统均通过 FUSE 对外提供 POSIX 接口。
本课题选择 FUSE 作为 RucksFS 的对外接入层，正是看重其"内核桥接 + 用户态逻辑"的分层模式——
可以用 Rust 这类对内核开发不友好的语言实现核心逻辑，
同时复用 Linux VFS 已经标准化的 POSIX 语义框架。

\subsubsection{架构与工作原理}

FUSE 的工作可以拆成两段来看：一段是\textbf{内核把文件系统调用路由到 \texttt{/dev/fuse}}，
另一段是\textbf{用户态守护进程从 \texttt{/dev/fuse} 读取请求并写回响应}。
图~\ref{fig:fuse_architecture} 以 \texttt{/dev/fuse} 为分界，将这两段并列呈现。

\begin{figure}[htbp]
    \centering
    \begin{tikzpicture}[
        font=\small,
        box/.style={draw, rounded corners=2pt, minimum width=3.2cm, minimum height=0.8cm,
                     align=center},
        kern/.style={box, fill=black!6},
        user/.style={box, fill=white},
        bridge/.style={draw, rounded corners=2pt, minimum width=2.0cm, minimum height=0.8cm,
                       align=center, fill=black!12, font=\small\ttfamily},
        arr/.style={->, >=Stealth, thick},
        lbl/.style={font=\scriptsize, text=black!65},
    ]

    % ======== 左侧通路：内核路由 ========
    \node[user] (app)     at (-4.2, 4.6) {应用程序};
    \node[kern] (vfs)     at (-4.2, 3.3) {Linux VFS};
    \node[kern] (fusemod) at (-4.2, 2.0) {FUSE 内核模块};

    \draw[arr] (app)     -- node[right, lbl] {syscall} (vfs);
    \draw[arr] (vfs)     -- node[right, lbl] {转发} (fusemod);

    % ======== 中央：/dev/fuse 桥接 ========
    \node[bridge] (devfuse) at (0, 2.0) {/dev/fuse};

    % ======== 右侧通路：用户态处理 ========
    \node[user] (daemon)   at (4.2, 2.0) {FUSE 守护进程};
    \node[user] (callback) at (4.2, 3.3) {回调派发};
    \node[user] (logic)    at (4.2, 4.6) {业务逻辑 / 存储};

    \draw[arr] (daemon)   -- node[right, lbl] {解析 opcode} (callback);
    \draw[arr] (callback) -- node[right, lbl] {调用} (logic);

    % ======== 跨 /dev/fuse 的双向箭头 ========
    \draw[arr] (fusemod) -- node[above, lbl] {请求} (devfuse);
    \draw[arr] (devfuse) -- node[above, lbl] {\texttt{read}} (daemon);

    % 响应返回通路（弧线，避免与下行箭头重叠）
    \draw[arr, bend left=28] (logic.west) to
        node[above, lbl] {响应} (devfuse.north);
    \draw[arr, bend left=28] (devfuse.north) to
        node[above, lbl, pos=0.7] {\texttt{write}} (fusemod.north);

    % ======== 图例：通过填充色区分用户/内核态 ========
    \node[kern, minimum width=0.8cm, minimum height=0.4cm] at (-2.6, 0.5) {};
    \node[lbl, anchor=west] at (-2.1, 0.5) {内核态组件};
    \node[user, minimum width=0.8cm, minimum height=0.4cm] at ( 1.2, 0.5) {};
    \node[lbl, anchor=west] at ( 1.7, 0.5) {用户态组件};

    \end{tikzpicture}
    \caption{FUSE 的内核路由与用户态处理。左侧是应用的文件系统调用在内核中的路径：
             系统调用进入 VFS，VFS 识别出挂载点类型为 \texttt{fuse} 后转交 FUSE 内核模块，
             最终把请求写入 \texttt{/dev/fuse} 字符设备队列并阻塞等待响应。
             右侧是用户态守护进程的处理路径：通过 \texttt{read(/dev/fuse)} 取出请求，
             解析 opcode 后派发到相应回调执行业务逻辑，再通过 \texttt{write(/dev/fuse)}
             把响应写回内核。两侧只通过 \texttt{/dev/fuse} 一条字节流通道交互。}
    \label{fig:fuse_architecture}
\end{figure}

\textbf{内核侧的路由}。当应用调用 \texttt{open("/mnt/fs/foo")}、\texttt{stat()}、\texttt{read()}
等系统调用时，请求首先进入 Linux VFS。VFS 查看路径所属挂载点的超级块，
发现该挂载点的文件系统类型是 \texttt{fuse}，于是把请求转交给 FUSE 内核模块（\texttt{fuse.ko}），
而不是像 ext4 那样直接调用内核内部的文件系统方法。FUSE 模块把 VFS 的结构化调用
编码为 FUSE 协议的字节序列——一个携带 opcode、请求 ID 和参数的帧——
然后写入与挂载点对应的 \texttt{/dev/fuse} 字符设备队列，
并让当前调用线程阻塞在该请求 ID 上，等待响应。

\textbf{用户态侧的处理}。用户态守护进程在挂载时已经打开了 \texttt{/dev/fuse}，
并在其上循环 \texttt{read}。一旦有请求到达，\texttt{read} 返回字节流，
守护进程按 FUSE 协议头解析出 opcode 与参数，把请求派发到开发者注册的回调上
（例如 \texttt{FUSE\_LOOKUP} 派发到 \texttt{lookup()}、\texttt{FUSE\_CREATE} 派发到 \texttt{create()}）。
回调内部可以自由调用任何用户态代码——读写 RocksDB、发起 gRPC、访问远端对象存储——
返回结果后，守护进程把结果按协议格式封装为响应帧，\texttt{write(/dev/fuse)} 送回内核；
FUSE 内核模块以请求 ID 为索引唤醒原先阻塞的应用线程，把返回值填回系统调用的返回路径。
如此循环，守护进程就像一个事件驱动的服务器：\texttt{/dev/fuse} 是它的 socket，
FUSE 协议是它的 RPC，回调函数是它的 handler。
开发者要做的就是把这些回调填上业务代码——下一节讨论具体的回调形式
以及 fuse3 如何让回调模型与 Rust 的异步运行时协作。

\subsubsection{开发接口与回调模型}

FUSE 框架向开发者暴露一组回调函数接口，每个回调对应一个 VFS 操作。
开发者实现所需的回调，注册到 FUSE 框架后，框架就会从 \texttt{/dev/fuse}
读取请求、分发到对应的回调、将返回值写回内核。
表~\ref{tab:fuse_callbacks} 列出了与本课题直接相关的核心回调。

\begin{table}[htbp]
    \centering
    \caption{FUSE 核心回调接口（本课题涉及部分）}
    \label{tab:fuse_callbacks}
    \begin{tabular}{llp{6.5cm}}
    \toprule
    \textbf{回调} & \textbf{对应 VFS 操作} & \textbf{语义} \\
    \midrule
    \texttt{lookup}   & 路径解析 & 给定父 inode 和文件名，返回子 inode 的属性 \\
    \texttt{getattr}  & \texttt{stat()} & 返回 inode 的元数据（mode、size、mtime 等） \\
    \texttt{setattr}  & \texttt{chmod/chown/utimes/truncate} & 修改 inode 属性 \\
    \texttt{create}   & \texttt{open(O\_CREAT)} & 在父目录下创建文件并返回文件句柄 \\
    \texttt{open}     & \texttt{open()} & 为已有文件分配 file handle \\
    \texttt{release}  & \texttt{close()} & 释放 file handle 及相关资源 \\
    \texttt{mkdir}    & \texttt{mkdir()} & 在父目录下创建子目录 \\
    \texttt{unlink}   & \texttt{unlink()} & 从父目录中删除文件的目录项 \\
    \texttt{rmdir}    & \texttt{rmdir()} & 删除空目录 \\
    \texttt{rename}   & \texttt{rename()} & 将目录项从源位置移动到目标位置 \\
    \texttt{readdir}  & \texttt{getdents()} & 列出目录下的所有条目 \\
    \texttt{read}     & \texttt{read()} & 读取文件数据 \\
    \texttt{write}    & \texttt{write()} & 写入文件数据 \\
    \texttt{fsync}    & \texttt{fsync()} & 将文件数据与元数据刷盘 \\
    \bottomrule
    \end{tabular}
\end{table}

原始 libfuse（C 语言）给回调提供的是阻塞同步接口：
守护进程主循环用一个线程对 \texttt{/dev/fuse} 做阻塞 \texttt{read}，
拿到请求后分派给线程池中的 worker 同步执行回调，
worker 拿到返回值后阻塞 \texttt{write} 写回响应。
这种"一请求一线程"的模型在元数据密集场景下会遇到两个问题：
一是回调中若要访问远端服务（如 RucksFS 的 gRPC 元数据调用），
线程会长时间阻塞在网络 I/O 上，占用线程池槽位；
二是线程数需要显著大于并发请求数才能避免排队，
而大量线程又会带来可观的内存和调度开销。

Rust 生态里的 \texttt{fuse3} crate 用异步运行时重写了这条主循环。
具体做法是把 \texttt{/dev/fuse} 的文件描述符注册到 tokio 的 \texttt{epoll} 反应器上，
开启 \texttt{EPOLLIN} 监听：内核一旦往设备队列写入请求，
\texttt{epoll\_wait} 就在 tokio 线程上返回可读事件，
fuse3 紧接着对 \texttt{/dev/fuse} 做非阻塞 \texttt{read}，
把取出的一批请求各自封装成独立的 \texttt{async} 任务（\texttt{tokio::task::spawn}），
这些任务就被派发到 tokio 的工作线程池。
每个任务调用对应的 \texttt{async fn} 回调——在 fuse3 中，
\texttt{Filesystem} trait 的 \texttt{lookup}、\texttt{create}、\texttt{getattr} 等方法
全部被声明为 \texttt{async fn}。
回调内部执行 \texttt{metadata\_client.create(...).await} 这类远程调用时，
底层的 \texttt{.await} 会把当前任务挂起，让出工作线程；
待 gRPC 响应到达、future 被唤醒后，任务回到工作线程完成剩余逻辑，
最后把响应帧 \texttt{write} 回 \texttt{/dev/fuse}。

这套"epoll + async 回调"模型的效果是：
数千个并发 FUSE 请求可以在少量工作线程上轻量地交错推进，
不再被阻塞式的网络 I/O 占用线程资源。
RucksFS 的 tokio 运行时只开启 8 个工作线程（等于 CPU 核数），
就足以承担第四章中 mdtest 在 $N=32$ 并发下的元数据负载。

\subsubsection{性能特征}

理解 FUSE 的性能开销，最直接的方式是把它和内核态文件系统的调用路径做对比。
对 ext4 这类内核态文件系统，一次 \texttt{stat()} 的路径是：
应用 $\to$ 内核态系统调用 $\to$ VFS $\to$ ext4 代码 $\to$ 返回用户态——
除了系统调用自身的一次 ring 0/3 切换外，整条路径都在内核内完成，
没有任何跨地址空间的数据拷贝。
而对 FUSE 文件系统，由于业务逻辑在另一个用户态进程里，
同一次 \texttt{stat()} 必须经过图~\ref{fig:fuse_architecture} 所示的完整往返：
应用 $\to$ 内核（路由）$\to$ \texttt{/dev/fuse}（请求入队）$\to$ 守护进程（处理）$\to$
\texttt{/dev/fuse}（响应入队）$\to$ 内核（匹配 ID）$\to$ 应用。
整个过程里发生\textbf{两次用户态/内核态上下文切换}
（应用线程睡眠再被唤醒、守护进程从 \texttt{epoll\_wait} 返回再回到 \texttt{write}），
以及\textbf{两次经 \texttt{/dev/fuse} 的数据拷贝}（请求帧和响应帧各一次）。
应用线程在整段过程中一直阻塞，真正的计算只发生在守护进程处理阶段——
剩下的时间基本都花在切换和拷贝上。

Vangoor 等人\cite{vangoor2017fuse}对 FUSE 相对内核态文件系统的开销做过系统测量：
在元数据操作（\texttt{getattr}、\texttt{lookup}）上，
FUSE 相比对应的内核态文件系统退化约 20\%--60\%——
因为这类操作本身执行时间极短（微秒级），固定的切换和拷贝开销占比极高；
数据操作（顺序读写大文件）退化较小，约 5\%--20\%，
因为此时磁盘或网络 I/O 时间远大于切换开销，分摊下来固定开销的影响被稀释。
本工作在第四章的分层开销实验中测到，
RucksFS 从进程内直接调用 RocksDB 的 172\,814\,ops/s，
到暴露为 FUSE 挂载点后的 2\,313\,ops/s，
FUSE 层本身只贡献了约 1.5$\times$ 的退化；
真正的大头（48.6$\times$）来自用户态的 gRPC 网络栈，
因为 RucksFS 的业务逻辑还涉及跨机器调用 MetadataServer。
换句话说，对本课题这种元数据密集、且元数据存储在远端服务的系统，
FUSE 的切换开销是不可消除但也不是主导因素的常数项——
优化的空间更多在 gRPC 栈和 RocksDB 访问上（第三章和第四章详述）。


\subsection{RocksDB 存储引擎}

文件系统的元数据——inode 属性、目录项、扩展属性等——天然适合用键值对的形式组织。
RocksDB\cite{dong2021rocksdb}是 Meta（原 Facebook）基于 Google
LevelDB\cite{google2011leveldb}开发的嵌入式键值存储引擎，
针对 SSD 和大规模写入场景做了大量优化。
MySQL 的 MyRocks 引擎、TiKV、CockroachDB 等生产系统均以 RocksDB 作为底层存储。

\subsubsection{LSM-tree 数据结构}

RocksDB 的底层数据结构为 LSM-tree（Log-Structured Merge-tree），
最早由 O'Neil 等人\cite{oneil1996lsmtree}于 1996 年提出。
其核心思想是将随机写转化为顺序写，以适应磁盘和 SSD 的 I/O 特性。
这一思想的渊源可以追溯到更早的日志结构文件系统（Log-Structured File System, LFS）\cite{rosenblum1992lfs}，
Rosenblum 和 Ousterhout 在 1991 年提出将所有写入——包括元数据——以日志追加的方式顺序写入磁盘，
LSM-tree 继承了这种"先写后整理"的策略并将其应用于索引结构。

写入路径如下：数据首先写入内存中的 MemTable（RocksDB 默认使用 SkipList 实现），
同时追加一条 WAL（Write-Ahead Log）记录以保证持久性。
当 MemTable 大小达到阈值（默认 64\,MB），它被转换为只读的 Immutable MemTable，
由后台线程 flush 到磁盘，生成一个有序的 SST（Sorted String Table）文件，
进入 Level-0 层。新的空 MemTable 随即可用于接收后续写入。
Level-0 的 SST 文件之间可能存在键范围重叠，当文件数量超过触发阈值
（默认 4 个）时，后台 Compaction 线程将 Level-0 文件与 Level-1 文件合并排序，
消除重叠并生成新的 SST 文件。类似的 Compaction 过程在更高层级之间逐层推进，
每一层的总容量以固定倍率（默认 10 倍）递增。

读取路径则相反：先查活跃 MemTable，再查 Immutable MemTable，
然后从 Level-0 开始逐层搜索 SST 文件，直到找到目标键或确认不存在。
为了加速点查询，每个 SST 文件内部维护索引块和可选的 Bloom 过滤器，
内存中还维护块缓存（Block Cache）存放频繁访问的数据块。

这种"写入友好"的结构使得 RocksDB 在写入密集场景下表现出色，
但也带来写放大（同一数据在 Compaction 过程中被反复读写）和
读放大（查询可能穿透多个层级）的问题。
RocksDB 提供了 Leveled、Universal、FIFO 三种 Compaction 策略，
允许使用者根据工作负载特征权衡写放大与空间利用率。


\begin{figure}[htbp]
    \centering
    \begin{tikzpicture}[
        sst/.style={draw, minimum width=0.9cm, minimum height=0.55cm, align=center, font=\scriptsize},
        lbl/.style={font=\footnotesize, text=black!60},
        arr/.style={->, >=stealth, thick, dashed, black!40}
    ]

    % MemTable
    \node[sst, fill=black!8, minimum width=2.2cm] (mem) at (0, 4) {MemTable};
    \node[lbl, anchor=east] at (-1.5, 4) {内存};

    % Level 0（允许重叠）
    \node[sst, fill=black!5] (l0a) at (-1, 2.8) {SST};
    \node[sst, fill=black!5] (l0b) at ( 0, 2.8) {SST};
    \node[sst, fill=black!5] (l0c) at ( 1, 2.8) {SST};
    \node[lbl, anchor=east] at (-1.5, 2.8) {L0};

    % Level 1
    \node[sst, fill=black!5] (l1a) at (-1.5, 1.7) {SST};
    \node[sst, fill=black!5] (l1b) at (-0.5, 1.7) {SST};
    \node[sst, fill=black!5] (l1c) at ( 0.5, 1.7) {SST};
    \node[sst, fill=black!5] (l1d) at ( 1.5, 1.7) {SST};
    \node[lbl, anchor=east] at (-1.5, 1.7) {L1};

    % Level 2
    \node[sst, fill=black!5] (l2a) at (-2, 0.6) {SST};
    \node[sst, fill=black!5] (l2b) at (-1, 0.6) {SST};
    \node[sst, fill=black!5] (l2c) at ( 0, 0.6) {SST};
    \node[sst, fill=black!5] (l2d) at ( 1, 0.6) {SST};
    \node[sst, fill=black!5] (l2e) at ( 2, 0.6) {SST};
    \node[lbl, anchor=east] at (-1.5, 0.6) {L2};

    % 箭头
    \draw[arr] (mem.south) -- ++(0, -0.45) node[right, lbl] {Flush};
    \draw[arr] (0, 2.45)   -- ++(0, -0.35) node[right, lbl] {Compaction};
    \draw[arr] (0, 1.35)   -- ++(0, -0.35) node[right, lbl] {Compaction};

    % 容量标注
    \node[lbl] at (3.2, 2.8) {可重叠};
    \node[lbl] at (3.2, 1.7) {有序不重叠};
    \node[lbl] at (3.2, 0.6) {容量 $\times 10$};

    \end{tikzpicture}
    \caption{LSM-tree 分层结构。写入先进入内存 MemTable，写满后 Flush 为 Level~0 的 SST 文件（允许键范围重叠），各层间通过 Compaction 合并排序，每层容量以 10 倍递增。}
    \label{fig:lsm_tree}
\end{figure}


\subsubsection{Column Family 机制}

Column Family（列族）是 RocksDB 在单个数据库实例内提供的逻辑分区机制，
设计上受到了 Bigtable\cite{chang2006bigtable}列族概念的启发。
每个 Column Family 拥有独立的 MemTable、SST 文件集合和 Compaction 配置，
但共享同一个 WAL 文件。一个 RocksDB 实例可以创建任意数量的 Column Family，
各自配置不同的压缩算法、block\_size、Bloom 过滤器参数和 Compaction 触发阈值。

这种隔离性在文件系统场景中有实际好处。inode 属性数据（几百字节一条）
以点查询为主，适合配置较小的 block\_size 来减少读放大，并启用全键 Bloom 过滤器。
目录项数据（父 inode 编号 + 文件名的映射）以前缀扫描为主——给定父目录 inode，
遍历其下所有子项——适合配置前缀提取器和前缀 Bloom 过滤器。
将两类数据放在不同的 Column Family 中，可以分别选择最优的配置，
避免单一配置下的折衷。

各 Column Family 独立 flush 和 Compaction 还有一个间接好处：
频繁更新的数据（如文件大小、修改时间戳）产生的 Compaction I/O
不会干扰相对稳定的目录结构数据。
同时，同一个 WriteBatch 可以包含对多个 Column Family 的操作，
并以原子方式提交——跨 Column Family 的原子性由共享 WAL 来保证。


\begin{figure}[htbp]
    \centering
    \begin{tikzpicture}[
        cf/.style={draw, rounded corners=2pt, minimum width=2.6cm, minimum height=1.8cm,
                    align=center, font=\small, fill=black!4},
        wal/.style={draw, rounded corners=2pt, minimum width=9.5cm, minimum height=0.6cm,
                    align=center, font=\small, fill=black!8},
        arr/.style={->, >=stealth, thick, black!40}
    ]

    % Column Families
    \node[cf] (cf1) at (0, 2.5)   {\textbf{inodes CF}\\[2pt]{\footnotesize MemTable}\\{\footnotesize SST files}\\{\footnotesize 全键 Bloom}};
    \node[cf] (cf2) at (3.5, 2.5) {\textbf{dir\_entries CF}\\[2pt]{\footnotesize MemTable}\\{\footnotesize SST files}\\{\footnotesize 前缀 Bloom}};
    \node[cf] (cf3) at (7, 2.5)   {\textbf{delta\_entries CF}\\[2pt]{\footnotesize MemTable}\\{\footnotesize SST files}\\{\footnotesize 前缀 Bloom}};

    % Shared WAL
    \node[wal] (wal) at (3.5, 0.4) {共享 Write-Ahead Log（WAL）};

    % Connections
    \draw[arr] (cf1.south) -- (cf1 |- wal.north);
    \draw[arr] (cf2.south) -- (cf2 |- wal.north);
    \draw[arr] (cf3.south) -- (cf3 |- wal.north);

    \end{tikzpicture}
    \caption{RocksDB 列族机制。每个列族拥有独立的 MemTable、SST 文件和 Bloom 过滤器配置，共享同一个 WAL 以保证跨列族原子写入。}
    \label{fig:column_family}
\end{figure}


\subsubsection{WriteBatch 与事务支持}

WriteBatch 是 RocksDB 的批量写入接口。
多个 \texttt{put()}、\texttt{delete()}、\texttt{merge()} 操作可以被收集到
同一个 WriteBatch 中，随后一次性提交。
从持久化角度看，整个 WriteBatch 只产生一次 WAL 写入和一次 \texttt{fsync}，
相比逐条提交将 I/O 次数降低了一到两个数量级。
从原子性角度看，WriteBatch 中的所有操作要么全部对读取者可见，
要么全部不可见，不存在部分提交的中间状态。

单纯的 WriteBatch 无法处理"读取—判断—修改"模式的操作，
因为在读取和写入之间，其他线程可能修改了相同的键。
RocksDB 的 TransactionDB 在 WriteBatch 基础上增加了行级锁机制。
打开数据库时指定 \texttt{TransactionDBOptions}，
即可使用 \texttt{begin\_transaction()} 创建事务对象。
事务内调用 \texttt{get\_for\_update(key)} 读取某个键时，
TransactionDB 对该键加排他锁，阻止其他事务修改同一键，
直到当前事务 \texttt{commit()} 或 \texttt{rollback()}。

TransactionDB 同时支持两种并发控制模式：
悲观模式（\texttt{TransactionDB}）在读取时即加锁；
乐观模式（\texttt{OptimisticTransactionDB}）在提交时才检测冲突。
文件系统元数据操作——特别是 \texttt{mkdir}、\texttt{rename}
等涉及父目录并发修改的操作——冲突概率相对较高，
悲观模式通过提前加锁避免了回滚的浪费。


\begin{figure}[htbp]
    \centering
    \begin{tikzpicture}[
        box/.style={draw, rounded corners=2pt, minimum width=2.8cm, minimum height=0.7cm,
                     align=center, font=\small, fill=black!4},
        arr/.style={->, >=stealth, thick},
        lbl/.style={font=\footnotesize, text=black!60}
    ]

    % 三个阶段
    \node[box] (collect)  at (0, 0)   {收集操作};
    \node[box] (wal)      at (4, 0)   {写入 WAL};
    \node[box] (memtable) at (8, 0)   {应用到 MemTable};

    \draw[arr] (collect)  -- (wal);
    \draw[arr] (wal)      -- (memtable);

    % 子标注
    \node[lbl, text width=2.5cm, align=center] at (0, -0.9)
        {Put A, Put B,\\Delete C, Merge D};
    \node[lbl, text width=2.5cm, align=center] at (4, -0.9)
        {单次 \texttt{fsync}\\持久化保证};
    \node[lbl, text width=2.5cm, align=center] at (8, -0.9)
        {所有操作\\原子可见};

    \end{tikzpicture}
    \caption{WriteBatch 执行流程。多个 put/delete 操作被收集到同一个批次中，通过一次 WAL 写入持久化，再原子地应用到 MemTable——全部可见或全部不可见。}
    \label{fig:write_batch}
\end{figure}


\subsubsection{Bloom 过滤器与前缀迭代器}

RocksDB 在每个 SST 文件中嵌入 Bloom 过滤器\cite{rocksdb2023wiki}，
用于在点查询时快速判断目标 key 是否可能存在于该文件中。
命中否定结果时可以跳过整个 SST 文件，避免无效的磁盘读取。
RucksFS 的 \texttt{inodes} 列族配置了全键 Bloom 过滤器（10\,bit/key），
用于加速 \texttt{getattr} 等点查操作。

除全键 Bloom 外，RocksDB 还提供前缀 Bloom 过滤器（Prefix Bloom Filter），
配合 \texttt{prefix\_extractor} 选项使用。
指定前缀提取函数后，同一前缀下的所有 key 共享一个 Bloom 条目；
配合 \texttt{prefix\_same\_as\_start} 选项，
迭代器在扫描完当前前缀后自动终止。

前缀迭代器直接决定了目录遍历的执行效率。
RucksFS 的 \texttt{dir\_entries} 列族将目录项键设计为"父 inode 编号（大端序 8 字节）+ 文件名"的复合格式，
前缀提取器设为 \texttt{FixedPrefix(9)}（1 字节类型前缀 + 8 字节 inode）。
\texttt{readdir} 操作通过前缀迭代器只扫描属于目标目录的条目，
时间复杂度与该目录的条目数成正比，而非与整个数据库规模相关。

\subsection{POSIX 文件系统语义}

POSIX\cite{ieee2017posix} 是本工作的语义合约。
RucksFS 虽然用 Rust 写、跑在用户态、元数据放在 RocksDB，
但对上层应用它必须"看起来像一个符合 POSIX 的文件系统"——否则依赖 POSIX 假设的工具
（cp、rsync、make、Git、sqlite\ldots）会在不起眼的细节上出错。
本节梳理与 RucksFS 设计直接相关的几组 POSIX 约束，
后续第三章的许多设计取舍都要回到这些语义要求上来理解。

\subsubsection{inode 模型与目录结构}

POSIX 标准定义了 Unix 系文件系统的基本抽象模型。
每个文件系统对象（普通文件、目录、符号链接、字符/块设备、FIFO、套接字等）由一个 inode（index node）唯一标识。
inode 编号在同一文件系统内全局唯一，存储了对象的全部元数据信息：
文件类型与权限模式（\texttt{mode}，16 位整数，其中高 4 位编码类型、低 12 位编码权限及 setuid/setgid/sticky）、
所有者和所属组（\texttt{uid}/\texttt{gid}，各 32 位）、
文件大小（\texttt{size}，64 位字节数，对目录特指占用字节，不是条目数）、
硬链接计数（\texttt{nlink}）、
访问时间（\texttt{atime}）、修改时间（\texttt{mtime}）、
状态变更时间（\texttt{ctime}，注意与 mtime 的区分——mtime 反映文件内容变化，
ctime 反映 inode 元数据本身变化，包括 chmod、chown、rename 等）。

目录是一种特殊文件，其数据内容是一组 (名称, inode 编号) 的映射关系，
即目录项（directory entry, dentry）。
同一个 inode 可以被多个目录项引用，这就是硬链接机制。
inode 的 \texttt{nlink} 字段记录指向它的硬链接数量。
当 \texttt{nlink} 降为 0 且没有进程持有该文件的打开文件描述符时，
文件系统回收该 inode 及其数据块。
对目录而言，\texttt{nlink} 有更特殊的含义：一个空目录的 \texttt{nlink} 初始为 2（
自身的目录项 + 自身下的 \texttt{.}），每增加一个子目录再 +1（子目录下的 \texttt{..} 指回来）。
这一看似冗余的规则是 POSIX 明文要求，工具如 \texttt{find} 会依赖 \texttt{nlink-2}
推算子目录数量来做剪枝，若实现不当会使这些工具效率显著下降。

根目录是整个文件系统命名空间的起点，通常使用 inode 编号 1 或 2（取决于具体实现）。
每个目录至少包含两个特殊条目：\texttt{.}（指向自身）和 \texttt{..}（指向父目录）。
路径解析从根目录开始，逐级在目录项中查找下一个分量的 inode 编号，
直到到达目标对象或遇到错误（如中间某级目录不存在）。
这个逐级查找过程在 FUSE 框架中对应 \texttt{lookup} 回调。
路径解析的每一步都要做权限检查——沿途每级目录都必须有执行位（\texttt{x}）——
因此深层路径的解析成本是 $O(\text{depth})$。
RucksFS 依赖内核 \texttt{default\_permissions} 完成这层检查，
用户态只负责返回正确的 \texttt{mode}、\texttt{uid}、\texttt{gid}。

\subsubsection{核心文件操作语义}

POSIX 标准对文件操作的行为有严格规定，
FUSE 文件系统必须正确实现这些语义才能被上层应用正常使用。
以下是几个关键操作的语义约束，也是 RucksFS 必须严格满足的合约。

\texttt{create} 在指定目录下创建一个新的普通文件。
系统需要分配一个未使用的 inode 编号、初始化 inode 元数据
（\texttt{nlink} 设为 1，\texttt{size} 设为 0，\texttt{mode} 按参数与 umask 的位运算结果设置），
并在父目录中插入一条目录项。
父目录的 \texttt{mtime} 和 \texttt{ctime} 需要同步更新。
如果同名文件已存在，返回 \texttt{EEXIST} 错误；
如果调用时带 \texttt{O\_EXCL} 但文件已存在，同样返回 \texttt{EEXIST}。
所有这些子步骤必须对其他进程表现为原子整体——不能观察到"目录项已出现但 inode 属性未写"之类的中间态。
RucksFS 通过单次 RocksDB \texttt{WriteBatch} 提交来保证这一原子性（详见第三章）。

\texttt{unlink} 从指定目录中删除一个目录项，
并将对应 inode 的 \texttt{nlink} 减 1。
如果 \texttt{nlink} 降至 0 \emph{且} 没有进程持有该文件的打开文件描述符，文件数据被释放；
\textbf{若仍有 fd 被打开，inode 和数据必须保留直到最后一个 fd 关闭}。
这一语义是 Unix 系统编程的常用技巧——
应用可以在 \texttt{open} 后立即 \texttt{unlink}，得到一个只被本进程可见的临时文件，
进程退出时系统自动回收。
为正确实现该语义，文件系统必须维护一个"待删除但还有 fd 引用"的 inode 集合；
RucksFS 对应的数据结构是 \texttt{pending\_deletes}（第三章第五节）。
\texttt{unlink} 只能用于非目录对象——
删除空目录需要使用 \texttt{rmdir}，
它会额外检查目录是否为空（仅包含 \texttt{.} 和 \texttt{..}）。
\texttt{unlink} 对一个目录调用时返回 \texttt{EISDIR}，
\texttt{rmdir} 对一个非空目录调用时返回 \texttt{ENOTEMPTY}。

\texttt{rename} 将一个目录项从源位置移动到目标位置。
POSIX 要求 \texttt{rename} 是原子操作：
其他进程要么看到旧状态，要么看到新状态，
不应观察到中间的不一致状态。
如果目标位置已有同名文件，\texttt{rename} 先隐式 \texttt{unlink} 目标文件，
这一步同样必须与移动操作原子一体。
跨目录的 \texttt{rename} 还需同时更新源父目录和目标父目录的元数据，
并检测是否构成目录循环——例如把 \texttt{/a} 重命名为 \texttt{/a/b}
应返回 \texttt{EINVAL}。
Linux 5.x 起还支持 \texttt{renameat2()} 系统调用的 \texttt{RENAME\_EXCHANGE} 模式，
原子交换两个路径所指向的 inode，这对 \texttt{sqlite}、\texttt{vim} 等需要"写入临时文件后原子替换"的工具有价值；
RucksFS 目前支持标准的 \texttt{rename}，\texttt{RENAME\_EXCHANGE} 留作后续工作。

\texttt{readdir} 遍历一个目录的所有条目，
返回 (名称, inode 编号, 文件类型) 三元组的列表。
POSIX 对遍历过程中的一致性保证较弱——
允许在遍历期间新增或删除的条目可见性不确定。
但有两项弱一致的不变式仍须保证：
一是同一 \texttt{opendir} 会话内的同一条目不会被读取两次；
二是在 \texttt{opendir} 前就存在且未被删除的条目一定会被读取一次。
RucksFS 的 \texttt{readdir} 实现通过在 RocksDB 上做前缀迭代（见 \ref{sec:bloom-prefix} 节），
并在迭代开始时建立快照来满足这两项不变式。

\texttt{open} 与 \texttt{release}（对应 \texttt{close}）之间维护了一份文件会话状态。
\texttt{open} 返回的文件描述符在同一进程内被 \texttt{dup} 后，
两个 fd 共享同一个会话（即 offset、访问模式等），而不是独立会话——
POSIX 严格区分"文件描述符"和"文件描述"：
前者是进程表中的条目，后者是内核中的会话结构。
FUSE 把这份会话暴露为 \texttt{open} 回调返回的 \texttt{fh}（file handle），
\texttt{release} 只在该会话的最后一个 fd 被关闭时才触发。
RucksFS 借此 handle 计数判断 \texttt{pending\_deletes} 中的 inode 能否真正回收。

\texttt{fsync} 和 \texttt{fdatasync} 要求将已写入的数据与元数据持久化到稳定存储，
系统调用返回时必须保证崩溃后数据仍在。
对 RucksFS 这一约束转化为：调用 MetadataServer 的 \texttt{fsync} 时必须触发
RocksDB 的 WAL \texttt{fsync} 和 DataServer 的 \texttt{fdatasync}，
两者任一失败都应返回 \texttt{EIO}。

\begin{table}[htbp]
    \centering
    \caption{与 RucksFS 实现直接相关的 POSIX 错误码约定}
    \label{tab:posix-errno}
    \begin{tabular}{llp{6.8cm}}
    \toprule
    \textbf{errno} & \textbf{常量} & \textbf{触发场景} \\
    \midrule
    2  & \texttt{ENOENT}    & 路径中某级不存在；\texttt{lookup}/\texttt{unlink} 找不到目标 \\
    13 & \texttt{EACCES}    & 权限不足；由 \texttt{default\_permissions} 代理返回 \\
    17 & \texttt{EEXIST}    & \texttt{create}/\texttt{mkdir} 时目标已存在 \\
    20 & \texttt{ENOTDIR}   & 以非目录作为路径中间分量 \\
    21 & \texttt{EISDIR}    & 对目录调用 \texttt{unlink} 或 \texttt{write} \\
    22 & \texttt{EINVAL}    & \texttt{rename} 形成循环、参数非法等 \\
    28 & \texttt{ENOSPC}    & 底层存储空间不足 \\
    30 & \texttt{EROFS}     & 对只读挂载执行写操作 \\
    39 & \texttt{ENOTEMPTY} & \texttt{rmdir} 时目录非空 \\
    11 & \texttt{EAGAIN}    & PCC 事务冲突，本文内部约定（见 \ref{sec:pcc}）\\
    \bottomrule
    \end{tabular}
\end{table}

表~\ref{tab:posix-errno} 汇总了 RucksFS 必须正确返回的几个核心错误码。
pjdfstest 测试套件中 8\,000 余条用例里的绝大多数，实质是在验证这些错误码
在各种边界输入下能否按 POSIX 规定触发。
第四章 \ref{sec:experiments} 节的 POSIX 合规性测试结果表明 RucksFS 在 8 类操作上 100\% 通过，
另外 4 类（link/symlink/mkfifo/mknod）由于对应的 inode 类型尚未实现而未通过。

\subsubsection{权限模型}

Unix 权限模型为每个 inode 关联三组 rwx（读/写/执行）权限位，
分别对应文件所有者、所属组和其他用户，编码在 \texttt{mode} 字段的低 9 位中；
加上 setuid、setgid、sticky 三位特殊权限，共 12 位。
权限检查发生在每次访问时，由内核 VFS 依据当前进程的 \texttt{euid}/\texttt{egid} 完成。
FUSE 框架提供了 \texttt{default\_permissions} 挂载选项，
启用后权限检查完全由内核 VFS 层代理完成，
用户态守护进程只需在 \texttt{getattr} 中正确返回 \texttt{mode}、\texttt{uid}、\texttt{gid} 即可；
对于沿路径遍历时的 \texttt{x} 权限检查、读目录时的 \texttt{r} 权限检查、
写文件时的 \texttt{w} 权限检查，用户态无需介入。
配合 \texttt{allow\_other} 选项，非挂载者的用户也可访问文件系统。
RucksFS 启用了这两个选项，将权限检查完全委托给内核，
既降低了实现复杂度，又保证了行为与内核原生文件系统一致。


\subsection{悲观并发控制}

文件系统的元数据操作天然面临并发竞争。多个进程可能同时在同一目录下创建文件、
删除文件或执行 rename。如果不加以控制，并发操作可能导致链接计数错误、
目录项残留或孤儿 inode 等一致性问题。

\subsubsection{PCC 基本原理}

悲观并发控制（Pessimistic Concurrency Control, PCC）的策略是"先加锁，再操作"。
事务在访问数据之前即获取排他锁，持有锁期间其他事务对同一数据的访问被阻塞，
直到锁释放。这种方式牺牲了一定的并行度，但换来了确定性的执行结果——
事务一旦开始，就不会因为冲突而被迫回滚。

与之对应的乐观并发控制（OCC）在事务执行期间不加锁，
仅在提交时检测冲突，若有冲突则回滚并重试。
OCC 在冲突率较低的场景下性能更好，但文件系统元数据操作（特别是同一目录下的并发创建和删除）
冲突概率较高，OCC 会导致大量无效回滚。
以 \texttt{mkdir} 为例：多个线程同时在同一父目录下创建子目录，
每次操作都需要检查子目录名是否已存在、插入目录项、更新父目录时间戳，
这些操作之间存在真实的写-写冲突。
PCC 通过在读取父目录 key 时就加锁，避免了这类无效计算。

\subsubsection{行级锁与事务隔离}

RocksDB TransactionDB 采用的锁粒度为单个 key——
可以视为关系数据库行级锁在键值存储中的等价物。
调用 \texttt{get\_for\_update(key)} 时，TransactionDB 对该 key 加排他锁。
如果另一个事务已持有同一 key 的锁，当前事务进入等待状态。
等待超时后（可配置，默认 1000\,ms），返回锁超时错误，调用方可选择重试或放弃。

锁表存储在内存中的哈希表里，不持久化到磁盘。
按 key 的哈希值分片（默认 16 个分片），减少了锁表本身的访问竞争。
当事务 \texttt{commit()} 或 \texttt{rollback()} 时，持有的所有锁被释放。

这种机制提供的隔离级别类似于 Read Committed：
事务读取到的数据是已提交的最新版本，不会看到其他未提交事务的中间结果。
结合排他锁，对于"读取—判断—修改"模式的操作，
从加锁到提交之间不会有其他事务修改同一 key。

在文件系统场景中，锁的粒度足够细。
一个 \texttt{create} 操作通常只需锁定父目录的目录项 key（检测同名、插入新条目）
和 inode 分配计数器 key（原子递增 inode 编号）。
如果两个进程同时在同一目录下创建不同名称的文件，
它们锁定的目录项 key 不同，可以并行执行。
只有名称相同时，后到的事务才会在 \texttt{get\_for\_update} 处阻塞。

\subsubsection{TOCTOU 竞态问题}

TOCTOU（Time-of-Check to Time-of-Use）是并发编程中一类常见的竞态条件。
程序先检查某个条件是否成立，再基于检查结果执行操作，
但在检查和操作之间存在时间窗口，其他线程可能在此窗口内改变了条件，
导致操作基于过期的信息执行。

以 \texttt{rmdir} 操作为例。按照 POSIX 语义，\texttt{rmdir} 只能删除空目录。
一个朴素的实现分两步执行：
第一步，遍历目标目录的条目列表，确认目录为空；
第二步，删除该目录的 inode 和父目录中的对应条目。
如果这两步不在同一个原子事务内完成，就存在如下竞态：
线程 A 在第一步确认目录为空后、第二步执行前，
线程 B 在该目录下创建了新文件。
随后线程 A 执行第二步删除目录——一个非空目录被错误地删除，
其中的文件变成了无法通过任何路径访问的孤儿 inode。

类似的 TOCTOU 风险存在于多个元数据操作中。
\texttt{rename} 在检查目标是否存在和执行移动之间有窗口；
\texttt{create} 在检查文件名不存在和插入目录项之间也有竞态；
\texttt{unlink} 在读取 \texttt{nlink} 和执行减一操作之间同样存在并发修改的可能。

解决 TOCTOU 的核心方法是将检查和操作纳入同一个原子事务。
在 RocksDB TransactionDB 的 PCC 模式下，
通过 \texttt{get\_for\_update} 在读取目录项时即加锁，
保证从检查到修改的整个过程中，其他事务无法插入或删除该目录的条目。
检查和操作之间的时间窗口被排他锁所覆盖，竞态条件随之消除。

回到 \texttt{rmdir} 的例子：事务开始后，
对目标目录的子条目前缀范围调用 \texttt{get\_for\_update}，确认为空且阻止并发插入，
再执行实际的删除操作，最后 \texttt{commit()} 原子提交。
如果在锁定过程中发现目录非空，事务直接回滚并返回 \texttt{ENOTEMPTY}。
整个过程中，检查和删除在同一个锁保护范围内完成。

\subsection{本章小结}

本章介绍了 RucksFS 所依赖的四项基础技术及其在本课题中的具体使用方式。
FUSE 提供了用户态文件系统的开发框架，开发者通过实现 lookup、create、readdir 等回调接口即可构建完整的文件系统，RucksFS 基于 fuse3 的异步模型开发。FUSE 的上下文切换开销（元数据操作退化 20\%--60\%）是可接受的常数因子，不影响本课题对存储层设计的验证。
RocksDB 的 LSM-tree 结构天然适合写入密集的元数据负载；Column Family 将不同类型的元数据隔离配置，WriteBatch 保证跨列族的原子写入，TransactionDB 提供行级锁的并发控制能力。Bloom 过滤器和前缀迭代器则分别加速了点查和目录遍历操作。
POSIX 语义定义了 create、unlink、rename、readdir 等操作的行为约束，这些约束直接决定了存储层的事务边界——rename 必须原子，unlink 必须正确维护链接计数。权限检查通过 \texttt{default\_permissions} 挂载选项委托给内核。
悲观并发控制通过在操作前加锁，将 TOCTOU 窗口归零。对于文件系统元数据这种冲突概率较高的场景，PCC 比乐观方案更经济。RucksFS 在此基础上实现了按 inode ID 排序加锁、冲突自动重试等具体策略。
